\documentclass[11pt]{article}
\usepackage[letterpaper,margin=0.85in]{geometry}\usepackage[T1]{fontenc}\usepackage[utf8]{inputenc}\usepackage{amsmath,amssymb,enumitem,booktabs,array,xcolor,fancyhdr,tabularx}
\setlist{itemsep=2pt,topsep=3pt}\pagestyle{fancy}\fancyhf{}\renewcommand{\headrulewidth}{0.3pt}
\fancyhead[L]{\small 3.4 Derivatives of Inverse Trigonometric Functions}\fancyhead[R]{\small Make-up key}\fancyfoot[C]{\small\thepage}
\definecolor{ans}{RGB}{28,125,74}\newcommand{\A}[1]{\textcolor{ans}{\textbf{#1}}}
\begin{document}

\section*{3.4 Derivatives of Inverse Trigonometric Functions}
\textbf{You missed this lesson. Copy this into your notes. That is all you need to do.}\par\medskip
\subsection*{Question of the day}
How fast does an angle change when its sine, or its tangent, changes?\par
\subsection*{Targets}\begin{itemize}\item I can say what arcsin, arccos and arctan mean.
\item I can find the derivative of an inverse trig function, with the chain rule.
\item I can find a slope or a tangent line for an inverse trig curve.\end{itemize}
\subsection*{Concepts}\begin{itemize}
\item \textbf{What arcsin means} (y = arcsin x  ?  sin y = x): arcsin x is the angle whose sine is x. sin\textsuperscript{-1} x means arcsin x, not 1/sin x.
\item \textbf{The three rules} (d/dx arcsin x = 1/$\surd$(1 -- x\textsuperscript{2})      d/dx arctan x = 1/(1 + x\textsuperscript{2})): arcsin and arccos have a square root. arctan has no root and a plus sign.
\item \textbf{The sign pair} (d/dx arccos x = --1/$\surd$(1 -- x\textsuperscript{2})): arccos has the opposite sign of arcsin. Everything else is the same.
\item \textbf{With the chain rule} (d/dx arcsin u = u? / $\surd$(1 -- u\textsuperscript{2})): Put u in for x, then multiply by u?. Square the whole inside: (3x)\textsuperscript{2} = 9x\textsuperscript{2}.
\item \textbf{Where they come from} (1 = cos y $\cdot$ dy/dx): Write x = sin y and differentiate. A right triangle with hypotenuse 1 turns cos y into $\surd$(1 -- x\textsuperscript{2}).
\end{itemize}
\subsection*{What the class predicted, and what is true}\begin{itemize}
\item Sine has a flat tangent at its peak. Mirror that piece across $y=x$. At $x=1$, the arcsin curve has a tangent that is... \A{vertical.} Rise and run swap in a mirror. A flat tangent becomes a vertical tangent, so the slope of arcsin grows without bound near $x=1$.
\item The graph of arctan levels off for large $x$. As $x$ gets larger, its slope gets... \A{closer to 0.} A curve that levels off has a tangent that gets flatter and flatter, so the slope heads toward 0.
\end{itemize}
\section*{Huddle notes (answers filled in)}
\subsection*{Definitions}\begin{itemize}
\item \textbf{arcsin x:} The \A{angle} whose sine is x. It is not 1/sin x.
\item \textbf{Inverse rule:} If y = arcsin x, then sin y = \A{x}. Differentiate both sides to find dy/dx.
\item \textbf{The three rules:} arcsin: 1/$\surd$(1 -- x\textsuperscript{2}). arccos: \A{--1}/$\surd$(1 -- x\textsuperscript{2}). arctan: 1/(1 + \A{x\textsuperscript{2}}).
\item \textbf{Chain rule:} Put u in for x, then multiply by \A{u?}. Square the \A{whole inside}.
\end{itemize}
\subsection*{Worked example}
Differentiate y = arcsin(3x).\\
\begin{itemize}
\item The inside is u = \A{3x}
\item Its derivative is u? = \A{3}
\item Square the whole inside: u\textsuperscript{2} = \A{9x\textsuperscript{2}}
\item So dy/dx = 3 / $\surd$(1 -- 9x\textsuperscript{2})
\end{itemize}
\subsection*{Label the triangle}
\begin{enumerate}
\item The longest side has length \A{1}.
\item The side opposite y has length \A{x}.
\item The third side has length \A{$\surd$(1 -- x\textsuperscript{2})}.
\end{enumerate}
\section*{One worked sample of each skill}
\subsection*{Differentiate}
\noindent\textbf{Level 1.} Differentiate.\\ \A{Answer: $\frac{-1}{\sqrt{1 - x^{2}}}$} \quad Use the arccos rule directly.\par\medskip
\noindent\textbf{Level 2.} Differentiate.\\ \A{Answer: $\frac{5}{\sqrt{1 - \left(5 x\right)^{2}}}$} \quad The inside is $u=5x$ and $u'=5$. Put them in the arcsin rule.\par\medskip
\noindent\textbf{Level 3.} Differentiate.\\ \A{Answer: $\frac{2 x}{1 + \left(x^{2}\right)^{2}}$} \quad The inside is $u=x^{2}$ and $u'=2x$. Put them in the arctan rule.\par\medskip
\noindent\textbf{Level 4.} Differentiate.\\ \A{Answer: $\frac{\frac{-1}{x^{2}}}{1 + \left(\frac{1}{x}\right)^{2}}$} \quad The inside is $u=\dfrac1x$ and $u'=-\dfrac1{x^2}$. Put them in the arctan rule.\par\medskip
\subsection*{Slope at a point}
\noindent\textbf{Level 1.} $y=\arctan\,x$. Find $\dfrac{dy}{dx}$ at $x=3$.\\ \A{Answer: 0.1} \quad $u(3)=3$ and $u'(3)=1$. Put them in the arctan rule: the slope is $0.1$.\par\medskip
\noindent\textbf{Level 2.} $y=\arcsin\left(5x\right)$. Find $\dfrac{dy}{dx}$ at $x=0.12$.\\ \A{Answer: 6.25} \quad $u(0.12)=0.6$ and $u'(0.12)=5$. Put them in the arcsin rule: the slope is $6.25$.\par\medskip
\noindent\textbf{Level 3.} $y=\arcsin\left(0.6e^{x}\right)$. Find $\dfrac{dy}{dx}$ at $x=0$.\\ \A{Answer: 0.75} \quad $u(0)=0.6$ and $u'(0)=0.6$. Put them in the arcsin rule: the slope is $0.75$.\par\medskip
\noindent\textbf{Level 4.} $y=\arcsin\left(0.6e^{x}\right)$. Find $\dfrac{dy}{dx}$ at $x=0$.\\ \A{Answer: 0.75} \quad $u(0)=0.6$ and $u'(0)=0.6$. Put them in the arcsin rule: the slope is $0.75$.\par\medskip
\subsection*{Spot the right rule}
\noindent\textbf{Level 1.} Which expression is $\dfrac{dy}{dx}$ for $y=\arccos\,x$?\\ (A) 1 / (1 + x\textsuperscript{2}) \quad (B) --1 / $\surd$(1 -- x) \quad (C) 1 / $\surd$(1 -- x\textsuperscript{2}) \quad \A{(D) --1 / $\surd$(1 -- x\textsuperscript{2})}\\ \A{Answer: D} \quad The inside is $u=x$, so $u'=1$ and $u^2=x^2$. The arccos rule gives --1 / $\surd$(1 -- x\textsuperscript{2}).\par\medskip
\noindent\textbf{Level 2.} Which expression is $\dfrac{dy}{dx}$ for $y=\arcsin\left(5x\right)$?\\ \A{(A) 5 / $\surd$(1 -- 25x\textsuperscript{2})} \quad (B) 5 / (1 + 25x\textsuperscript{2}) \quad (C) 5 / $\surd$(1 -- 5x) \quad (D) 1 / $\surd$(1 -- 25x\textsuperscript{2})\\ \A{Answer: A} \quad The inside is $u=5x$, so $u'=5$ and $u^2=25x^2$. The arcsin rule gives 5 / $\surd$(1 -- 25x\textsuperscript{2}).\par\medskip
\noindent\textbf{Level 3.} Which expression is $\dfrac{dy}{dx}$ for $y=\arctan\left(x^{2}\right)$?\\ (A) 1 / (1 + x?) \quad (B) 2x / (1 + x\textsuperscript{2}) \quad (C) 2x / $\surd$(1 -- x?) \quad \A{(D) 2x / (1 + x?)}\\ \A{Answer: D} \quad The inside is $u=x^{2}$, so $u'=2x$ and $u^2=x^4$. The arctan rule gives 2x / (1 + x?).\par\medskip
\noindent\textbf{Level 4.} Which expression is $\dfrac{dy}{dx}$ for $y=\arctan\left(e^{x}\right)$?\\ (A) e? / (1 + e?) \quad (B) e? / $\surd$(1 -- e\textsuperscript{2}?) \quad \A{(C) e? / (1 + e\textsuperscript{2}?)} \quad (D) 1 / (1 + e\textsuperscript{2}?)\\ \A{Answer: C} \quad The inside is $u=e^{x}$, so $u'=e^{x}$ and $u^2=e^{2x}$. The arctan rule gives e? / (1 + e\textsuperscript{2}?).\par\medskip
\subsection*{Tangent line}
\noindent\textbf{Level 1.} Write the tangent line to $y=\arctan\,x$ at $x=0$.\\ \A{Answer: $1 \left(x - 0\right) + 0$} \quad The point is $(0,0)$ and the slope is $1$. So $y=1(x-0)+0$.\par\medskip
\noindent\textbf{Level 2.} Write the tangent line to $y=\arctan\left(4x\right)$ at $x=0.25$.\\ \A{Answer: $2 \left(x - 0.25\right) + \frac{\pi}{4}$} \quad The point is $(0.25,\tfrac{\pi}{4})$ and the slope is $2$. So $y=2(x-0.25)+\tfrac{\pi}{4}$.\par\medskip
\noindent\textbf{Level 3.} Write the tangent line to $y=\arctan\left(\dfrac1x\right)$ at $x=1$.\\ \A{Answer: $\left(-0.5\right) \left(x - 1\right) + \frac{\pi}{4}$} \quad The point is $(1,\tfrac{\pi}{4})$ and the slope is $-0.5$. So $y=-0.5(x-1)+\tfrac{\pi}{4}$.\par\medskip
\noindent\textbf{Level 4.} Write the tangent line to $y=\arctan\left(\dfrac1x\right)$ at $x=1$.\\ \A{Answer: $\left(-0.5\right) \left(x - 1\right) + \frac{\pi}{4}$} \quad The point is $(1,\tfrac{\pi}{4})$ and the slope is $-0.5$. So $y=-0.5(x-1)+\tfrac{\pi}{4}$.\par\medskip
\end{document}